% !TITLE 过零丁洋
% !AUTHOR 文天祥
% !DYNASTY 两宋
% !GENRE 七言律诗
% !ORDER 18

\begin{poem}
辛苦遭逢起一经\textsuperscript{1}，干戈\textsuperscript{2}寥落四周星\textsuperscript{3}。
山河破碎风飘絮\textsuperscript{4}，身世浮沉雨打萍\textsuperscript{5}。
惶恐滩\textsuperscript{6}头说惶恐，零丁洋\textsuperscript{7}里叹零丁。
人生自古谁无死？留取丹心照汗青\textsuperscript{8}。
\end{poem}

\begin{annotations}
\notetext{1}{起一经：靠着精通一种经书起家。文天祥是进士出身，通过科举（考儒家经义）走上仕途。}
\notetext{2}{干戈：代指战争。南宋末年与元朝的战争已经持续了四年。}
\notetext{3}{四周星：四周年。岁星（木星）绕太阳一周约十二年，四周星就是四年。}
\notetext{4}{风飘絮：被风吹散的柳絮。山河像被风吹散的柳絮一样四分五裂。}
\notetext{5}{雨打萍：被雨敲打的浮萍。自己的一生像无根的浮萍一样漂泊不定。}
\notetext{6}{惶恐滩：赣江十八滩中最险的一滩，在今江西万安县。文天祥兵败后曾经过此地。“惶恐”一词双关——既是地名，也是心情。}
\notetext{7}{零丁洋：珠江口的一片海域，在今广东中山市南。零丁洋对文天祥来说是人生最后一段路——他在这里被元军俘虏。“零丁”一词双关——既是地名，也指孤苦零丁的处境。}
\notetext{8}{汗青：古代在竹简上写字之前，先用火烤竹子使它出水（如出汗），称为“汗青”，后来代指史册。我要留下一颗赤诚之心照耀在史册上——这句成了中国人最有名的就义宣言。}
\end{annotations}

\begin{appreciation}
《过零丁洋》是南宋最后一位丞相文天祥的绝命诗，也是中国文学史上最有名的就义诗之一。

文天祥的起点很平常：苦读经书，二十岁考中状元，靠精通一经走上仕途。可这个"起一经"的读书人，偏偏要独自面对一个时代的崩塌。"干戈寥落四周星"——和元军的仗打了四年，越打越没有希望，他还在打。"山河破碎风飘絮，身世浮沉雨打萍"是全诗最悲凉的一联：山河像被风吹散的柳絮，飘到哪里算哪里；自己像雨里的浮萍，沉下去浮上来都不由自己做主。国和己两个比喻摆在一起，你就明白他站在怎样的绝境里了。

"惶恐滩头说惶恐，零丁洋里叹零丁"妙得像天意。惶恐滩是赣江上真实的地名，他兵败路过；零丁洋是珠江口的海域，他在这里被俘。地名里恰好嵌着自己的处境，仿佛命运替他写好了词。读到这里，人是往下沉的。

最后两句却猛地提了起来："人生自古谁无死？留取丹心照汗青。"人都是要死的，这不是新道理；新的在后半句——死法可以自己选。汗青就是史册。古人写竹简前先烤竹子，竹面像出汗，所以叫汗青。他要做的事很朴素：让自己配得上将要写进史册的那个名字。

这句话让我们想起《离骚》。屈原说"虽九死其犹未悔"，文天祥说"人生自古谁无死"——一个在汨罗江边，一个在零丁洋上，隔着一千七百年，说的是同一件事：有些东西比命重要。写完这首诗，南宋灭亡，他被押到大都，拒绝一切劝降，就义时衣带里藏着一句遗言："孔曰成仁，孟曰取义"。

哥哥选这首诗，其实想跟你讲一句反着的话。英雄用死成全自己，我们普通人没有这种机会，也没有这种必要。咱们家的"史册"就是这本小书，等你看完它，翻回这一页，还记得今天有一个不怕死的人——那你的人生任务就完成了：好好活着，在要紧的事上不亏心。这已经很难了。
\end{appreciation}
